\documentclass[10pt,a4paper]{amsart}
\usepackage[margin=19mm]{geometry}
\usepackage{amsmath,amssymb,mathtools}
\usepackage[scr=rsfs,cal=euler]{mathalfa}
\usepackage{xfrac}
\usepackage[unicode,hidelinks,bookmarks=false]{hyperref}
\newcommand{\ie}{i.e.\ }
\newcommand{\cf}{cf.\ }
\newcommand{\resp}{respectively\ }
\newcommand{\Subsection}{\subsection}
\providecommand{\nobreakdash}{\nobreak\hbox{-}}
\setlength{\emergencystretch}{2em}
\multlinegap0pt
\usepackage{amssymb}
\usepackage[shortlabels]{enumitem}
\usepackage{xcolor}
\newtheorem{theorem}{Theorem}
\newcommand{\E}{\mathbb{E}}
\renewcommand{\P}{\mathbb{P}}
\title{The Blown Lead Paradox\\[0.25em]\large A Pathwise Calibration Benchmark for Win Probability Forecasts}
\author{Jonathan Pipping-Gam\'on, Abraham J. Wyner}
\date{Source: arXiv:2601.18774v4 (CC BY 4.0)}
\begin{document}
\maketitle

\noindent\emph{Source and licence.} The Blown Lead Paradox\\[0.25em]\large A Pathwise Calibration Benchmark for Win Probability Forecasts. Version arXiv:2601.18774v4 (CC BY 4.0), distributed by arXiv under the Creative Commons Attribution 4.0 International licence, http://creativecommons.org/licenses/by/4.0/, as recorded in the arXiv licence metadata for this exact version and shown on the abstract page https://arxiv.org/abs/2601.18774v4. Full licence notice: \texttt{licenses/arxiv-2601.18774-CC-BY-4.0.txt}.\par\smallskip

\noindent\textit{Editorial scope.} Counted unit: the article's theorem thm:supp-discrete-unconditional (source0.tex lines 777--786). The article's complete original proof of it is delivered with it (source0.tex lines 788--812, one \texttt{proof} environment of 25 lines). No mathematics is written, repaired or re-derived: the statement and the proof are the article's own lines, in the article's own notation, and no retained line is substituted or re-flowed.  No further source-local context is needed: every cross-reference of every retained line resolves inside the two delivered spans.  Nothing was omitted from any retained span, so the package carries no omission record.  No retained line contains a citation, so the wrapper carries no bibliography and no bibliography entry.  No retained line contains a figure environment, an image-inclusion command or drawing code, so no image is delivered and the package declares no asset dependency.  Editorial formatting only, in addition: an AMS \texttt{amsart} preamble that restates the article's own declarations and macros used by the retained lines, namely the environments \texttt{theorem} on separate counters, together with the standard packages those lines need.  The retained environments print in this document as Theorem 1 in order of appearance, whereas the article prints the same items with the numbering of its own full document.  The complete native source of the article as distributed in the arXiv e-print for this exact version is bundled unchanged as \texttt{sources/193416/source0.tex}.
\begin{theorem}[Discrete-time unconditional bound] \label{thm:supp-discrete-unconditional}
Let $M_N=\max_{0\le k\le N} p_k$. For $x\in[p_0,1)$,
\[
\P(M_N<x)\;\ge\;1-\frac{p_0}{x},
\]
and hence $F_{M_N}(x)=\P(M_N\le x)\ge 1-p_0/x$.
The strict-threshold bound is an equality if $p_{\tau_x}=x$ a.s.\ on $\{\tau_x\le N\}$; equality for $F_{M_N}(x)$ additionally requires $\P(M_N=x)=0$.
For $x<p_0$, $F_{M_N}(x)=0$.
Moreover, $M_N$ has an atom at $1$ with $\P(M_N=1)=p_0$.
\end{theorem}

\begin{proof}
If $x<p_0$ then $M_N\ge p_0>x$ a.s., hence $F_{M_N}(x)=0$.

Fix $x\in[p_0,1)$. By optional stopping for the bounded martingale $(p_k)$ at $\tau_x\wedge N$,
\[
\E[p_{\tau_x\wedge N}] = p_0.
\]
Decompose $p_{\tau_x\wedge N}=p_{\tau_x}\mathbf 1\{\tau_x\le N\}+p_N\mathbf 1\{\tau_x>N\}$.
On $\{\tau_x\le N\}$ we have $p_{\tau_x}\ge x$.
Therefore,
\[
p_0
= \E[p_{\tau_x}\mathbf 1\{\tau_x\le N\}] + \E[p_N\mathbf 1\{\tau_x>N\}]
\;\ge\; x\,\P(\tau_x\le N).
\]
Since $\{\tau_x\le N\}=\{M_N\ge x\}$, this yields
\[
\P(M_N\ge x)\le \frac{p_0}{x}\qquad\Rightarrow\qquad
\P(M_N<x)\ge 1-\frac{p_0}{x}.
\]
Equality in the strict-threshold bound holds when $p_{\tau_x}=x$ a.s.\ on $\{\tau_x\le N\}$.
Because $F_{M_N}(x)=\P(M_N<x)+\P(M_N=x)$, the corresponding CDF equality also requires no atom at $x$.

Finally, $M_N=1$ iff $p_N=1$ iff $Y=1$, so $\P(M_N=1)=p_0$.
\end{proof}

\end{document}
